% !TeX program = xelatex
% UTF-8; compile twice with XeLaTeX. No external images or data files needed.
\documentclass[a4paper,landscape,10pt]{article}
\usepackage[margin=14mm,headheight=14pt,headsep=5mm,footskip=7mm]{geometry}
\usepackage{fontspec}
\setmainfont{DejaVu Sans}
\newfontfamily\ipafont{DejaVu Sans}
\usepackage[table]{xcolor}
\usepackage{array,longtable,booktabs,fancyhdr}
\definecolor{ink}{HTML}{18384D}
\definecolor{band}{HTML}{E7EFF4}
\pagestyle{fancy}\fancyhf{}
\fancyhead[L]{\small\color{ink}Parlez Vous QSO?! · Po polsku / Polish}
\fancyhead[R]{\small Englische Phrase · Polskie tłumaczenie · IPA}
\fancyfoot[L]{\scriptsize Vorlage: Polish v 1.0 · Credits: Chris -- SP3ELD}
\fancyfoot[R]{\small\thepage}
\setlength{\parindent}{0pt}
\setlength{\tabcolsep}{3mm}
\renewcommand{\arraystretch}{1.25}
\newcommand{\ipa}[1]{{\ipafont[#1]}}
\begin{document}
{\LARGE\bfseries\color{ink}Parlez Vous QSO?!}\hfill{\large Po polsku / Polish}

\vspace{2mm}
{\small QSO-Sprachführer mit polnischer IPA-Lautschrift. Die Lautschrift bezieht sich auf die polnische Spalte.}

\vspace{2mm}
{\footnotesize \textbf{Leseschlüssel:} \ipa{ˈ} Hauptbetonung, \ipa{ˌ} Nebenbetonung;
\ipa{ɨ} polnisches y; \ipa{w} entspricht ł, \ipa{v} entspricht w;
\ipa{ɕ, ʑ, tɕ, dʑ} stehen für ś, ź, ć, dź; \ipa{ʂ, ʐ, tʂ, dʐ} für sz, ż/rz, cz, dż;
\ipa{ɛ̃, ɔ̃} markieren Nasalvokale.}

\vspace{1mm}
\fontsize{9}{11.8}\selectfont
\begin{longtable}{@{}>{\raggedright\arraybackslash}p{13mm}>{\raggedright\arraybackslash}p{73mm}>{\raggedright\arraybackslash}p{75mm}>{\raggedright\arraybackslash}p{90mm}@{}}
\toprule
\textbf{Ref.} & \textbf{English phrase} & \textbf{Po polsku} & \textbf{Lautschrift (IPA)}\\
\midrule
\endfirsthead
\toprule
\textbf{Ref.} & \textbf{English phrase} & \textbf{Po polsku} & \textbf{Lautschrift (IPA)}\\
\midrule
\endhead
\midrule
\multicolumn{4}{r}{\scriptsize Fortsetzung auf der nächsten Seite}\\
\endfoot
\bottomrule
\endlastfoot
\rowcolor{band} 1.0 & \textbf{Opening Greetings} & \multicolumn{2}{l}{\textbf{Powitanie}} \\*
1.1 & I am very pleased to meet you. & Miło mi cię spotkać. & \ipa{ˈmiwɔ mi tɕɛ ˈspɔtkatɕ} \\
1.2 & Thank you for your call & Dziękuję za twoje zawołanie. & \ipa{dʑɛŋˈkujɛ za ˈtfɔjɛ zavɔˈwaɲɛ} \\
1.3 & Good morning & Dzień dobry & \ipa{dʑɛɲ ˈdɔbrɨ} \\
1.31 & Good day & Dzień dobry & \ipa{dʑɛɲ ˈdɔbrɨ} \\
1.32 & Good afternoon & Dzień dobry & \ipa{dʑɛɲ ˈdɔbrɨ} \\
1.33 & Good evening & Dobry wieczór & \ipa{ˈdɔbrɨ ˈvjɛtʂur} \\
1.4 & I know only a little Polish but I enjoy QSOs in Polish. & Mówię słabo po polsku, ale chciałbym spróbować QSO w tym języku. & \ipa{ˈmuvjɛ ˈswabɔ pɔ ˈpɔlsku ˈalɛ ˈxtɕawbɨm spruˈbɔvatɕ ku ɛs ɔ f tɨm jɛ̃ˈzɨku} \\
\rowcolor{band} 2.0 & \textbf{Operator} & \multicolumn{2}{l}{\textbf{Operator}} \\*
2.1 & My name is Peter. & Mam na imię Piotr. & \ipa{mam na ˈimjɛ pjɔtr} \\
2.2 & I have been licensed since 1960. & Licencję mam od roku 1960. & \ipa{liˈtsɛntsje mam ɔt ˈrɔku ˈtɨɕɔnts dʑɛˈvjɛɲtɕsɛt ʂɛɕtɕdʑɛˈɕɔntɛɡɔ} \\
2.3 & I have been licensed for 43 years. & Licencję mam od 43 lat. & \ipa{liˈtsɛntsje mam ɔt tʂtɛrˈdʑɛstu ˈtʂɛx lat} \\
2.4 & I live in Sittingbourne. & Mieszkam w Sittingbourne. & \ipa{ˈmjɛʂkam f ˈsɪtɪŋbɔːn} \\
2.5 & Sittingbourne is a village. & Sittingbourne to wioska. & \ipa{ˈsɪtɪŋbɔːn tɔ ˈvjɔska} \\
2.51 & Sittingbourne is a town. & Sittingbourne to miasteczko. & \ipa{ˈsɪtɪŋbɔːn tɔ mjaˈstɛtʂkɔ} \\
2.52 & Sittingbourne is a city. & Sittingbourne to miasto. & \ipa{ˈsɪtɪŋbɔːn tɔ ˈmjastɔ} \\
2.6 & Sittingbourne is between London and Dover. & Sittingbourne leży między Londynem a Dover. & \ipa{ˈsɪtɪŋbɔːn ˈlɛʐɨ ˈmjɛndzɨ lɔnˈdɨnɛm a ˈdəʊvə} \\
2.7 & I live in the South of England. & Mieszkam w południowej Anglii. & \ipa{ˈmjɛʂkam f pɔwudˈɲɔvɛj ˈaŋɡlji} \\
2.71 & I live in the North of England. & Mieszkam w północnej Anglii. & \ipa{ˈmjɛʂkam f puwˈnɔtsnɛj ˈaŋɡlji} \\
2.72 & I live in the East of England. & Mieszkam we wschodniej Anglii. & \ipa{ˈmjɛʂkam vɛ ˈfsxɔdɲɛj ˈaŋɡlji} \\
2.73 & I live in West of England. & Mieszkam w zachodniej Anglii. & \ipa{ˈmjɛʂkam v zaˈxɔdɲɛj ˈaŋɡlji} \\
2.74 & I live in South East England. & Mieszkam w południowo-wschodniej Anglii. & \ipa{ˈmjɛʂkam f pɔwudˌɲɔvɔ ˈfsxɔdɲɛj ˈaŋɡlji} \\
2.75 & I live in South West England. & Mieszkam w południowo-zachodniej Anglii. & \ipa{ˈmjɛʂkam f pɔwudˌɲɔvɔ zaˈxɔdɲɛj ˈaŋɡlji} \\
2.76 & I live in North East England. & Mieszkam w północno-wschodniej Anglii. & \ipa{ˈmjɛʂkam f puwˌnɔtsnɔ ˈfsxɔdɲɛj ˈaŋɡlji} \\
2.77 & I live in North West England. & Mieszkam w północno-zachodniej Anglii. & \ipa{ˈmjɛʂkam f puwˌnɔtsnɔ zaˈxɔdɲɛj ˈaŋɡlji} \\
2.78 & I live in central England. & Mieszkam w środkowej Anglii. & \ipa{ˈmjɛʂkam f ɕrɔtˈkɔvɛj ˈaŋɡlji} \\
2.8 & I live in the country. & Mieszkam na wsi. & \ipa{ˈmjɛʂkam na fɕi} \\
2.81 & I live by the sea. & Mieszkam nad morzem. & \ipa{ˈmjɛʂkam nad ˈmɔʐɛm} \\
2.82 & I live near mountains. & Mieszkam w górach. & \ipa{ˈmjɛʂkam v ˈɡurax} \\
2.9 & Sittingbourne is 50Km from London. & Sittingbourne leży 50 km od Londynu. & \ipa{ˈsɪtɪŋbɔːn ˈlɛʐɨ pjɛɲˈdʑɛɕɔnt kilɔˈmɛtruf ɔd lɔnˈdɨnu} \\
2.11 & I am retired. & Jestem emerytem. & \ipa{ˈjɛstɛm ɛmɛˈrɨtɛm} \\
2.12 & My home call sign is M0ECS. & Mój prywatny znak wywoławczy to M0ECS. & \ipa{muj prɨˈvatnɨ znak vɨvɔˈwaftʂɨ tɔ ɛm ˈzɛrɔ ɛ tsɛ ɛs} \\
\rowcolor{band} 3.0 & \textbf{Station} & \multicolumn{2}{l}{\textbf{Stacja}} \\*
3.1 & My radio is an Icom 706. & Moje radio to Icom 706. & \ipa{ˈmɔjɛ ˈradjɔ tɔ ˈaɪkɒm ˈɕɛdɛmsɛt ʂɛɕtɕ} \\
3.2 & My station is all homebrew (homemade). & Moja stacja w całości jest typu homemade. & \ipa{ˈmɔja ˈstatsja f tsaˈwɔɕtɕi jɛst ˈtɨpu ˌhəʊmˈmeɪd} \\
3.3 & I am running 5 Watts. & Nadaję z mocą 5 watów. & \ipa{naˈdajɛ z ˈmɔtsɔ̃ ˈpjɛɲtɕu ˈvatuf} \\
3.4 & My antenna is a long wire. & Moja antena to long wire. & \ipa{ˈmɔja anˈtɛna tɔ ˈlɒŋ waɪə} \\
3.41 & My antenna is a dipole. & Moja antena to dipol. & \ipa{ˈmɔja anˈtɛna tɔ ˈdipɔl} \\
3.42 & My antenna is a G5RV. & Moja antena to G5RV. & \ipa{ˈmɔja anˈtɛna tɔ ɡjɛ pjɛɲtɕ ɛr faw} \\
3.43 & My antenna is a doublet. & Moja antena to doublet. & \ipa{ˈmɔja anˈtɛna tɔ ˈdʌblət} \\
3.44 & My antenna is a delta loop. & Moja antena to delta. & \ipa{ˈmɔja anˈtɛna tɔ ˈdɛlta} \\
3.45 & My antenna is a vertical. & Moja antena to wertikal. & \ipa{ˈmɔja anˈtɛna tɔ vɛrˈtikal} \\
3.46 & My antenna is a log periodic. & Moja antena jest logarytmiczno-periodyczna. & \ipa{ˈmɔja anˈtɛna jɛst lɔɡarɨtˌmitʂnɔ pɛrjɔˈdɨtʂna} \\
3.47 & My antenna is a Yagi beam. & Moja antena to beam. & \ipa{ˈmɔja anˈtɛna tɔ biːm} \\
3.48 & My antenna is a 5 element beam. & Moja antena to beam 5-elementowy. & \ipa{ˈmɔja anˈtɛna tɔ biːm ˌpjɛɲtɕɔɛlɛmɛnˈtɔvɨ} \\
\rowcolor{band} 4.0 & \textbf{Conditions} & \multicolumn{2}{l}{\textbf{Warunki}} \\*
4.1 & Thank you for my 599 report. & Dziękuję za mój raport 599. & \ipa{dʑɛŋˈkujɛ za muj ˈrapɔrt pjɛɲtɕ ˈdʑɛvjɛɲtɕ ˈdʑɛvjɛɲtɕ} \\
4.2 & The band conditions are good. & Warunki na paśmie są dobre. & \ipa{vaˈruŋki na ˈpaɕmjɛ sɔ̃ ˈdɔbrɛ} \\
4.3 & The band conditions are bad. & Warunki na paśmie są złe. & \ipa{vaˈruŋki na ˈpaɕmjɛ sɔ̃ zwɛ} \\
4.4 & There is strong interference here. & Mam tutaj duże zakłócenia. & \ipa{mam ˈtutaj ˈduʐɛ zakwuˈtsɛɲa} \\
4.5 & Can you change to the 80m band? & Możesz przejść na pasmo 80 m? & \ipa{ˈmɔʐɛʂ pʂɛjɕtɕ na ˈpasmɔ ɔɕɛmˈdʑɛɕɔnt ˈmɛtruf} \\
4.6 & Please change frequency to 3.558 MHz. & Przejdź na częstotliwość 3,558 MHz. & \ipa{pʂɛjtɕ na tʂɛ̃stɔˈtlivɔɕtɕ tʂɨ ˈpʂɛtɕinɛk pjɛɲtɕ pjɛɲtɕ ˈɔɕɛm mɛɡaˈxɛrtsuf} \\
\rowcolor{band} 5.0 & \textbf{General Queries and Comments} & \multicolumn{2}{l}{\textbf{Ogólne pytania i komentarze}} \\*
5.1 & What does <\#\#\#\#\#> mean please? & Co znaczy <\#\#\#\#\#> ? & \ipa{tsɔ ˈznatʂɨ …} \\
5.2 & Please wait a moment, I need to check my Polish dictionary. & Zaczekaj proszę chwilę, muszę sprawdzić w słowniku polskim. & \ipa{zaˈtʂɛkaj ˈprɔʂɛ ˈxfʲilɛ ˈmuʂɛ ˈspravdʑitɕ f swɔvˈɲiku ˈpɔlskim} \\
5.3 & How do you say, <\#\#\#\#\#> in Polish? & Jak się mówi <\#\#\#\#\#> po polsku? & \ipa{jak ɕɛ ˈmuvi … pɔ ˈpɔlsku} \\
5.5 & My favourite modes of operation are CW / RTTY / PSK / MT63 / Olivia / Hellschreiber / Digital Modes / and SSB & Najchętniej nadaję na telegrafii / RTTY / PSK / MT63 / Olivia / Hellschreiber / emisjami cyfrowymi / na SSB. & \ipa{najˈxɛntɲɛj naˈdajɛ na tɛlɛˈɡrafji / ɛr tɛ tɛ ˈiɡrɛk / pɛ ɛs ka / ɛm tɛ ʂɛɕtɕˈdʑɛɕɔnt tʂɨ / ɔˈlivja / ˈhɛlˌʃʁaɪbɐ / ɛmiˈsjami tsɨfrɔˈvɨmi / na ɛs ɛs bɛ} \\
5.6 & My best DX is Australia. & Moim największym DX-em jest Australia. & \ipa{ˈmɔim najˈvjɛŋkʂɨm dɛ ˈiksɛm jɛst awsˈtralja} \\
5.7 & I enjoy working contests. & Lubię pracować w zawodach. & \ipa{ˈlubjɛ praˈtsɔvatɕ v zaˈvɔdax} \\
5.8 & I like working QRP. & Lubię nadawać QRP. & \ipa{ˈlubjɛ naˈdavatɕ ku ɛr pɛ} \\
\rowcolor{band} 6.0 & \textbf{The Weather} & \multicolumn{2}{l}{\textbf{Pogoda}} \\*
6.1 & It is sunny. & Jest słonecznie. & \ipa{jɛst swɔˈnɛtʂɲɛ} \\
6.2 & It is windy. & Jest wietrznie. & \ipa{jɛst ˈvjɛtʂɲɛ} \\
6.3 & It is raining. & Pada deszcz. & \ipa{ˈpada dɛʂtʂ} \\
6.4 & It is snowing. & Pada śnieg. & \ipa{ˈpada ɕɲɛk} \\
6.5 & It is very hot. & Jest bardzo gorąco. & \ipa{jɛst ˈbardzɔ ɡɔˈrɔntsɔ} \\
6.6 & It is very cold. & Jest bardzo zimno. & \ipa{jɛst ˈbardzɔ ˈʑimnɔ} \\
6.7 & The temperature is 8C. & Temperatura wynosi 8 stopni Celsjusza. & \ipa{tɛmpɛraˈtura vɨˈnɔɕi ˈɔɕɛm ˈstɔpɲi tsɛlˈsjuʂa} \\
6.8 & It is overcast. & Niebo jest zachmurzone. & \ipa{ˈɲɛbɔ jɛst zaxmuˈʐɔnɛ} \\
6.9 & It is foggy. & Jest mgliście. & \ipa{jɛst ˈmɡliɕtɕɛ} \\
6.11 & It has been raining all day. & Pada cały dzień. & \ipa{ˈpada ˈtsawɨ dʑɛɲ} \\
6.12 & There is 20cm of snow here. & Leży 20 cm śniegu. & \ipa{ˈlɛʐɨ ˈdvadʑɛɕtɕa tsɛntɨˈmɛtruf ˈɕɲɛɡu} \\
\rowcolor{band} 8.0 & \textbf{QSL} & \multicolumn{2}{l}{\textbf{Karta QSL}} \\*
8.1 & My QSL card is sure via the bureau & Kartę wyślę na pewno przez biuro. & \ipa{ˈkartɛ ˈvɨɕlɛ na ˈpɛvnɔ pʂɛz ˈbjurɔ} \\
8.2 & QSL via eQSL & Potwierdzę QSO kartą eQSL. & \ipa{pɔtˈfjɛrdzɛ ku ɛs ɔ ˈkartɔ̃ i ku ɛs ɛl} \\
8.3 & No eQSL please & Nie chcę eQSL. & \ipa{ɲɛ xtsɛ i ku ɛs ɛl} \\
8.4 & I will send you my card direct. & Wyślę ci moją kartę QSL direktem. & \ipa{ˈvɨɕlɛ tɕi ˈmɔjɔ̃ ˈkartɛ ku ɛs ɛl diˈrɛktɛm} \\
8.5 & What is your address? & Jaki jest twój adres? & \ipa{ˈjaki jɛst tfuj ˈadrɛs} \\
8.6 & Please send your QSL card direct. & Wyślij twą kartę QSL direktem. & \ipa{ˈvɨɕlij tfɔ̃ ˈkartɛ ku ɛs ɛl diˈrɛktɛm} \\
8.7 & My address is... & Mój adres jest następujący... & \ipa{muj ˈadrɛs jɛst nastɛmpuˈjɔntsɨ …} \\
8.8 & My address is shown on QSL.com / EHam.net & Mój adres jest dostępny na QSL.com / na eHam.net. & \ipa{muj ˈadrɛs jɛst dɔˈstɛmpnɨ na ku ɛs ɛl ˈkrɔpka kɔm / na i hæm ˈkrɔpka nɛt} \\
8.9 & Sorry I do not send QSL cards. & Przykro mi, ale nie wysyłam kart QSL. & \ipa{ˈpʂɨkrɔ mi ˈalɛ ɲɛ vɨˈsɨwam kart ku ɛs ɛl} \\
8.11 & Please send QSL card via G1JQH. & Wyślij proszę kartę via G1JQH. & \ipa{ˈvɨɕlij ˈprɔʂɛ ˈkartɛ ˈvja ɡjɛ ˈjɛdɛn jɔt ku xa} \\
8.12 & My WACRAL / GQRP / DXCC / FISTS number is 1234. & Mój numer WACRAL / GQRP / DXCC / FISTS jest 1234. & \ipa{muj ˈnumɛr ˈpɔdvujnɛ vu a tsɛ ɛr ɛl / ɡjɛ ku ɛr pɛ / dɛ iks tsɛ tsɛ / fɪsts jɛst ˈtɨɕɔnts ˈdvjɛɕtɕɛ ˈtʂɨdʑɛɕtɕi ˈtʂtɛrɨ} \\
\rowcolor{band} 9.0 & \textbf{Final Greetings} & \multicolumn{2}{l}{\textbf{Pożegnanie}} \\*
9.1 & Thank you. & Dziękuję. & \ipa{dʑɛŋˈkujɛ} \\
9.2 & Thank you for a very enjoyable contact. & Dziękuję ci za miłą łączność. & \ipa{dʑɛŋˈkujɛ tɕi za ˈmiwɔ̃ ˈwɔntʂnɔɕtɕ} \\
9.3 & I enjoyed our chat. & Jestem zadowolony z naszego chatu. & \ipa{ˈjɛstɛm zadɔvɔˈlɔnɨ z naˈʂɛɡɔ ˈtʂatu} \\
9.4 & I enjoyed our QSO in Polish. & Jestem zadowolony z naszej łączności po polsku. & \ipa{ˈjɛstɛm zadɔvɔˈlɔnɨ z ˈnaʂɛj wɔntʂˈnɔɕtɕi pɔ ˈpɔlsku} \\
9.5 & It has been great fun. & To była wielka przyjemność. & \ipa{tɔ ˈbɨwa ˈvjɛlka pʂɨˈjɛmnɔɕtɕ} \\
9.6 & If you would like a regular contact with me, you can contact me via email. & Jeśli chcesz się kontaktować ze mną regularnie, skontaktuj się ze mną e-mailem. & \ipa{ˈjɛɕli xtsɛʂ ɕɛ kɔntakˈtɔvatɕ zɛ mnɔ̃ rɛɡuˈlarɲɛ skɔnˈtaktuj ɕɛ zɛ mnɔ̃ iˈmɛjlɛm} \\
9.7 & My email address is... & Mój adres e-mailowy jest... & \ipa{muj ˈadrɛs imɛjˈlɔvɨ jɛst …} \\
9.8 & My email address is shown on QSL.com / EHam.net & Mój adres jest dostępny na QSL.com / na eHam.net. & \ipa{muj ˈadrɛs jɛst dɔˈstɛmpnɨ na ku ɛs ɛl ˈkrɔpka kɔm / na i hæm ˈkrɔpka nɛt} \\
9.9 & I look forward to meeting you again. & Mam nadzieję na następne spotkanie w przyszłości. & \ipa{mam naˈdʑɛjɛ na naˈstɛmpnɛ spɔtˈkaɲɛ f pʂɨˈʂwɔɕtɕi} \\
9.11 & My best wishes to you and your family. & Najlepsze życzenia dla ciebie i twojej rodziny. & \ipa{najˈlɛpʂɛ ʐɨˈtʂɛɲa dla ˈtɕɛbjɛ i ˈtfɔjɛj rɔˈdʑinɨ} \\
9.12 & Best wishes to your wife. & Najlepsze życzenia dla twej żony. & \ipa{najˈlɛpʂɛ ʐɨˈtʂɛɲa dla tfɛj ˈʐɔnɨ} \\
9.121 & Best wishes to your husband. & Najlepsze życzenia dla twego męża. & \ipa{najˈlɛpʂɛ ʐɨˈtʂɛɲa dla ˈtfɛɡɔ ˈmɛ̃ʐa} \\
9.13 & Good luck. & Życzę szczęścia! & \ipa{ˈʐɨtʂɛ ˈʂtʂɛ̃ɕtɕa} \\
9.14 & Good health. & Życzę zdrowia! & \ipa{ˈʐɨtʂɛ ˈzdrɔvja} \\
9.15 & Goodbye & Do widzenia! & \ipa{dɔ viˈdʑɛɲa} \\
9.16 & I wish you an enjoyable afternoon. & Miłego popołudnia. & \ipa{miˈwɛɡɔ pɔpɔˈwudɲa} \\
9.161 & I wish you an enjoyable evening. & Miłego wieczoru. & \ipa{miˈwɛɡɔ vjɛˈtʂɔru} \\
9.162 & I wish you an enjoyable weekend. & Miłego weekendu. & \ipa{miˈwɛɡɔ viˈkɛndu} \\
9.17 & Enjoy your holiday. & Przyjemnych wakacji! & \ipa{pʂɨˈjɛmnɨx vaˈkatsji} \\
9.18 & Happy New Year. & Szczęśliwego Nowego Roku! & \ipa{ʂtʂɛ̃ɕˈlivɛɡɔ nɔˈvɛɡɔ ˈrɔku} \\
\rowcolor{band} 10.0 & \textbf{Country} & \multicolumn{2}{l}{\textbf{Kraj}} \\*
10.1 & Albania & Albania & \ipa{alˈbaɲja} \\
10.2 & Austria & Austria & \ipa{ˈawstrja} \\
10.3 & Belgium & Belgia & \ipa{ˈbɛlɡja} \\
10.4 & Belorussia & Białoruś & \ipa{bjaˈwɔruɕ} \\
10.5 & Bosnia \& Herzegovina & Bośnia i Hercegowina & \ipa{ˈbɔɕɲa i xɛrtsɛɡɔˈvina} \\
10.6 & Bulgaria & Bułgaria & \ipa{buwˈɡarja} \\
10.7 & Croatia & Chorwacja & \ipa{xɔrˈvatsja} \\
10.8 & Czech Republic & Czechy & \ipa{ˈtʂɛxɨ} \\
10.9 & Denmark & Dania & \ipa{ˈdaɲja} \\
10.11 & Eire & Irlandia & \ipa{irˈlandja} \\
10.12 & England & Anglia & \ipa{ˈaŋɡlja} \\
10.13 & Estonia & Estonia & \ipa{ɛsˈtɔɲja} \\
10.14 & Finland & Finlandia & \ipa{finˈlandja} \\
10.15 & France & Francja & \ipa{ˈfrantsja} \\
10.16 & Germany & Niemcy & \ipa{ˈɲɛmtsɨ} \\
10.17 & Gibraltar & Gibraltar & \ipa{dʑiˈbraltar} \\
10.18 & Greece & Grecja & \ipa{ˈɡrɛtsja} \\
10.19 & Holland & Holandia & \ipa{xɔˈlandja} \\
10.21 & Hungary & Węgry & \ipa{ˈvɛŋɡrɨ} \\
10.22 & Iceland & Islandia & \ipa{isˈlandja} \\
10.23 & Northern Ireland & Irlandia Północna & \ipa{irˈlandja puwˈnɔtsna} \\
10.24 & Isle of Man & Wyspa Man & \ipa{ˈvɨspa man} \\
10.25 & Italy & Włochy & \ipa{ˈvwɔxɨ} \\
10.26 & Latvia & Łotwa & \ipa{ˈwɔtfa} \\
10.27 & Liechtenstein & Liechtenstein & \ipa{ˈliçtn̩ʃtaɪn} \\
10.28 & Lithuania & Litwa & \ipa{ˈlitfa} \\
10.29 & Luxembourg & Luksemburg & \ipa{lukˈsɛmburk} \\
10.31 & Macedonia & Macedonia & \ipa{matʂɛˈdɔɲja} \\
10.32 & Moldavia & Mołdawia & \ipa{mɔwˈdavja} \\
10.33 & Netherlands & Holandia & \ipa{xɔˈlandja} \\
10.34 & Norway & Norwegia & \ipa{nɔrˈvɛɡja} \\
10.35 & Poland & Polska & \ipa{ˈpɔlska} \\
10.36 & Portugal & Portugalia & \ipa{pɔrtuˈɡalja} \\
10.37 & Romania & Rumunia & \ipa{ruˈmuɲja} \\
10.38 & Russian Federation & Federacja Rosyjska & \ipa{fɛdɛˈratsja rɔˈsɨjska} \\
10.39 & Scotland & Szkocja & \ipa{ˈʂkɔtsja} \\
10.41 & Slovakia & Słowacja & \ipa{swɔˈvatsja} \\
10.42 & Slovenia & Słowenia & \ipa{swɔˈvɛɲja} \\
10.43 & Spain & Hiszpania & \ipa{xiʂˈpaɲja} \\
10.44 & Svalbard & Svalbard & \ipa{ˈsvalbart} \\
10.45 & Sweden & Szwecja & \ipa{ˈʂfɛtsja} \\
10.46 & Switzerland & Szwajcaria & \ipa{ʂfajˈtsarja} \\
10.47 & Turkey & Turcja & \ipa{ˈturtsja} \\
10.48 & Ukraine & Ukraina & \ipa{uˈkraina} \\
10.49 & Wales & Walia & \ipa{ˈvalja} \\
\end{longtable}
\normalsize
\section*{Hinweise zur Übertragung und Aussprache}
\fontsize{8.5}{10.5}\selectfont\raggedright
\textbf{Vorlage und Umfang.} Grundlage ist „Polish v 1.0.doc“; Credit: Chris -- SP3ELD.
Alle 142 Phrasen und Ländernamen sind enthalten. Die ursprüngliche Gliederung mit ihren Nummernsprüngen
und dem fehlenden Abschnitt 7 bleibt erhalten. Zwei vollständig leere Zusatzspalten der Vorlage wurden
entfernt; statt ihrer steht hier die IPA-Spalte.

\smallskip
\textbf{Dokumentierte Korrekturen.} „w wschodniej“ wurde zu „we wschodniej“ (2.72),
„Lichtenstein“ zu „Liechtenstein“ (10.27), „Luksemurg“ zu „Luksemburg“ (10.29) und
„Swalbard“ zu „Svalbard“ (10.44). Die Ländernamen entsprechen damit den Schreibweisen in polnischen
amtlichen Verzeichnissen. Mehrfache Leerzeichen wurden vereinheitlicht. Alle übrigen polnischen
Formulierungen bleiben wie in der Vorlage; die Lautschrift folgt dem korrigierten Text.

\smallskip
\textbf{Inhaltliche Eigenheiten.} Peter wird in 2.1 als Piotr übersetzt. In 1.4 steht sinngemäß
„Ich spreche schlecht Polnisch, möchte aber ein QSO in dieser Sprache versuchen“; diese Formulierung
bleibt erhalten. „near mountains“ wird in 2.82 mit „w górach“ wiedergegeben. Bei 3.47 fehlt Yagi in der
polnischen Spalte; dort steht nur „beam“. In 9.8 nennt der polnische Satz die Adresse ohne den ausdrücklichen
Zusatz „E-Mail“. Holland und Netherlands stehen beide als „Holandia“ (10.19 und 10.33).
Historische Ländernamen, Webadressen und englische Fachwörter der Vorlage wurden nicht modernisiert.
„chciałbym“, „emerytem“ und „zadowolony“ sind männliche Sprecherformen; für eine Sprecherin passen
„chciałabym“, „emerytką“ und „zadowolona“.

\smallskip
\textbf{Transkriptionskonvention.} Die IPA-Spalte ist eine breite, manuell erstellte Aussprachehilfe
für Standardpolnisch, keine enge Transkription einer Aufnahme. Hauptbetonung ist mit \ipa{ˈ}, ausgewählte
Nebenbetonung mit \ipa{ˌ} markiert. Polnische Wörter sind überwiegend auf der vorletzten Silbe betont;
bei Zahlen und längeren Wörtern können Ausnahmen auftreten. Polnisch hat keine bedeutungsunterscheidende Vokallänge; \ipa{ː} steht hier nur in fremdsprachigen Formen.

\smallskip
\textbf{Besondere Laute.} y ist \ipa{ɨ}, i ist \ipa{i}; ł ist \ipa{w} wie englisches w,
w ist \ipa{v} wie deutsches w. Die weichen Zischlaute \ipa{ɕ, ʑ} und Affrikaten \ipa{tɕ, dʑ}
unterscheiden sich von den härteren Lauten \ipa{ʂ, ʐ, tʂ, dʐ}. Affrikaten sind platzsparend ohne Bindebögen
geschrieben. j bezeichnet den Gleitlaut; vor i können Lippen- und andere Konsonanten zusätzlich palatalisiert
sein. Auslautverhärtung und Angleichungen sind berücksichtigt; feinere Details nur teilweise.

\smallskip
\textbf{Nasalvokale.} ą und ę haben je nach Umgebung verschiedene Aussprachen. Vor Verschlusslauten
und Affrikaten stehen häufig Vokal plus Nasalkonsonant, beispielsweise \ipa{ɛŋ} in „dziękuję“ und
\ipa{ɔntʂ} in „łączność“. Vor Frikativen und bei auslautendem ą wird hier breit mit \ipa{ɛ̃, ɔ̃}
transkribiert; die natürliche Aussprache enthält oft einen nasalen Gleitlaut. Auslautendes ę erscheint
in der häufigen Aussprache ohne deutliche Nasalierung, etwa „dziękuję“ als \ipa{dʑɛŋˈkujɛ}.


\smallskip
\textbf{Namen, Zahlen und Kürzel.} Sittingbourne und Dover sind englisch transkribiert;
Icom verwendet eine englische Aussprachevariante. Homemade, Long Wire, Doublet und Beam sind ebenfalls
englisch, Hellschreiber und Liechtenstein deutsch wiedergegeben. E-Mail und seine polnischen Endungen
werden polnisch angepasst gelesen. Rufzeichen und die meisten Kürzel stehen mit polnischen Buchstabennamen,
nicht im internationalen Buchstabieralphabet: Q „ku“, W „podwójne wu“, J „jot“, V hier „fał“.
FISTS wird als Wort gelesen; das e in eQSL und eHam als englisches \ipa{i}.
Zahlen stehen mit den zur Phrase passenden Formen: 43 nach „od“ als „czterdziestu trzech“,
5 watów als „pięciu watów“, 5-elementowy als „pięcioelementowy“. 706 wird als „siedemset sześć“,
599 als „pięć dziewięć dziewięć“, 1234 als „tysiąc dwieście trzydzieści cztery“ und
3,558 MHz als „trzy przecinek pięć pięć osiem megaherców“ gelesen.
Schrägstriche trennen Alternativen beziehungsweise Listenelemente; Platzhalter \texttt{<\#\#\#\#\#>}
erscheinen in der Lautschrift als Auslassung und werden nicht mitgesprochen.

\smallskip
\textbf{Fachquellen.} Katarzyna Dziubalska-Kołaczyk und Bogdan Walczak: „Polish“ (2011),
Universitätsrepositorium der Adam-Mickiewicz-Universität, Abschnitt 1.3.1 zur Phonologie;
Wiktor Jassem: „Polish“, Journal of the International Phonetic Association 33(1), 2003, S. 103--107.
Länderschreibweisen: polnisches Außenministerium / amtliche geografische Namensverzeichnisse.
Diese Quellen stützen die allgemeinen Konventionen; die einzelnen Phrasen wurden hier manuell transkribiert.

\smallskip
\textbf{LaTeX.} UTF-8-Datei; zweimal mit XeLaTeX kompilieren. Benötigt werden DejaVu Sans und die
üblichen LaTeX-Pakete der Präambel. IPA ist direkt als Unicode eingegeben; keine externen Dateien erforderlich.
\end{document}
